\documentclass[11pt]{article}

\usepackage{fullpage}
\usepackage{amsmath}
\usepackage{amssymb}
\usepackage{amsfonts}
\usepackage{mathtools}
\usepackage{enumitem}
\usepackage[colorlinks,linkcolor=magenta,citecolor=blue]{hyperref}

\newcommand{\Rbar}{\bar R}
\newcommand{\Om}{\Omega}
\newcommand{\E}{\mathbb E}
\newcommand{\tr}{\operatorname{tr}}
\newcommand{\diag}{\operatorname{diag}}
\newcommand{\Normal}{\mathcal N}
\newcommand{\Unif}{\operatorname{Unif}}
\newcommand{\IW}{\mathcal{IW}}
\newcommand{\Gam}{\operatorname{Gamma}}

\title{SER-IW and the Shared-\texorpdfstring{\(\omega\)}{omega} Approximation}
\author{}
\date{}

\begin{document}
\maketitle

\section{Goal}

The purpose of this note is to separate two objects:
\begin{enumerate}[leftmargin=*]
\item the more principled single-effect inverse-Wishart model, called SER-IW
here;
\item a tractable approximation that leads to a fast shared-\(\omega\) SER
variant.
\end{enumerate}
The shared-\(\omega\) approximation is not designed to make the SER PIP exactly
ideal. Instead, the goal is to make the approximation explicit, fast, and
amenable to later bounding of the SNP-specific PIP distortion.

\section{SER-IW Model}

Let \(x\in\mathbb R^J\) be the vector of marginal summary statistics,
\(\Rbar\) be the reference LD matrix, and \(N\) be the GWAS sample size. The
single-effect inverse-Wishart model is
\[
R\sim \IW(\nu_0\Rbar,\nu_0+J+1),
\qquad
\gamma\sim\Unif\{1,\ldots,J\},
\qquad
\beta\sim\Normal(0,\sigma_0^2),
\]
and
\[
x\mid R,\beta,\gamma=j
\sim
\Normal(\beta R_{\cdot j},R/N).
\]
This is the original natural model in which the same unknown in-sample LD
matrix \(R\) appears in both the mean and the noise covariance.

\subsection{Block decomposition}

For a fixed active SNP \(j\), write
\[
d_j=R_{jj},
\qquad
z_j=\frac{R_{-j,j}}{d_j},
\qquad
S_j=R_{-j,-j}-d_jz_jz_j^\top.
\]
Then
\[
R=
\begin{pmatrix}
d_j & d_jz_j^\top\\
d_jz_j & S_j+d_jz_jz_j^\top
\end{pmatrix}.
\]
Let
\[
\mu_j=\Rbar_{-j,j},
\qquad
\bar S_j=\Rbar_{-j,-j}-\mu_j\mu_j^\top.
\]
The inverse-Wishart block identities give the working representation
\[
d_j\sim \operatorname{IG}\left(\frac{\nu_0+2}{2},\frac{\nu_0}{2}\right),
\qquad
z_j\mid S_j
\sim
\Normal\left(\mu_j,\frac{S_j}{\nu_0}\right),
\]
and
\[
S_j\sim \IW(\nu_0\bar S_j,\nu_0+J+1),
\]
using the same convention as in the earlier notes.

\subsection{Useful likelihood factorization}

Under the SER-IW likelihood,
\[
x_j\mid d_j,\beta,\gamma=j
\sim
\Normal(\beta d_j,d_j/N),
\]
and the Gaussian conditioning formula gives
\[
x_{-j}\mid x_j,z_j,S_j,\gamma=j
\sim
\Normal(x_jz_j,S_j/N).
\]
The important point is that the conditional distribution of \(x_{-j}\) given
\(x_j\), \(z_j\), and \(S_j\) does not involve \(\beta\) or \(d_j\).

If we fix \(d_j\equiv1\), as in the model-B simplification, this becomes
\[
x_j\mid \beta,\gamma=j
\sim
\Normal(\beta,1/N),
\qquad
x_{-j}\mid x_j,z_j,S_j,\gamma=j
\sim
\Normal(x_jz_j,S_j/N).
\]
This factorization is the source of the ideal-PIP behavior in the
single-effect model: the active coordinate carries the \(\beta\)-evidence,
whereas the nonactive coordinates learn the LD column.

\subsection{\texorpdfstring{Exact conditional marginal under the \(d_j=1\) simplification}{Exact conditional marginal under the d-j equals 1 simplification}}

Conditioning on \(S_j\), the prior and likelihood for \(z_j\) are conjugate:
\[
z_j\mid S_j
\sim
\Normal(\mu_j,S_j/\nu_0),
\qquad
x_{-j}\mid x_j,z_j,S_j
\sim
\Normal(x_jz_j,S_j/N).
\]
Integrating over \(z_j\) gives
\[
x_{-j}\mid x_j,S_j,\gamma=j
\sim
\Normal\left(
x_j\mu_j,
\left(\frac1N+\frac{x_j^2}{\nu_0}\right)S_j
\right).
\]
After integrating over \(S_j\), this becomes a multivariate \(t\)-type
conditional distribution centered at \(x_j\mu_j\). Equivalently, it can be
represented using a normal--Gamma mixture:
\[
x_{-j}\mid x_j,\omega_j,\gamma=j
\sim
\Normal\left(
x_j\mu_j,
\frac{1}{\omega_j}
\left(\frac1N+\frac{x_j^2}{\nu_0}\right)\bar S_j
\right),
\]
\[
\omega_j
\sim
\Gam\left(\frac{\nu_0+3}{2},\frac{\nu_0}{2}\right),
\]
up to the precise scale convention used for the multivariate \(t\). This
conditional representation is close to ideal, but it is awkward to extend to
IBSS because the conditional covariance depends on the observed active
coordinate \(x_j\) and only describes \(x_{-j}\mid x_j\), not a common
Gaussian residual model for the whole vector.

\section{Approximation Leading to Shared-\texorpdfstring{\(\omega\)}{omega}}

The shared-\(\omega\) approximation replaces the exact conditional
matrix-valued uncertainty by a scalar local scale. For each candidate active
SNP \(j\), introduce \(\omega_j\) and set
\[
\omega_j
\sim
\Gam(a_\omega,b_\omega),
\qquad
a_\omega=\frac{\nu_0+3}{2},
\qquad
b_\omega=\frac{\nu_0}{2}.
\]
The same \(\omega_j\) scales both the observation noise and the relaxed-column
prior:
\[
x\mid \beta,\gamma=j,z_j,\omega_j
\sim
\Normal_J\left(
\beta u_j(z_j),
\frac{\Rbar}{N\omega_j}
\right),
\]
\[
z_j\mid \omega_j
\sim
\Normal_{J-1}\left(
\mu_j,
\frac{\bar S_j}{\nu_0\omega_j}
\right).
\]
Here \(u_j(z_j)\) is the \(J\)-vector whose \(j\)-th entry is one and whose
remaining entries are \(z_j\).

This approximation makes three deliberate simplifications.

\begin{enumerate}[leftmargin=*]
\item The random Schur complement \(S_j\) is replaced by its reference value
\(\bar S_j\), up to the scalar scale \(1/\omega_j\).
\item The same scalar \(\omega_j\) controls uncertainty in both the noise and
the LD column. This is what keeps the posterior update for \(\omega_j\)
conjugate.
\item The conditional model \(x_{-j}\mid x_j\) is replaced by a full-vector
Gaussian likelihood with covariance \(\Rbar/(N\omega_j)\). This is less exact
for SER PIPs, but much easier to embed in an IBSS residual algorithm.
\end{enumerate}

\subsection*{My attempt to justify the approximation}
Write the model as 
\begin{equation}
x = \beta \begin{pmatrix} 1 \\ z_1 \end{pmatrix} + z_0,
\end{equation}
where
\begin{equation}\label{eq:generative-z0-z1-by-R}
\begin{cases}
z_1 \mid S_1 \sim \Normal(\mu_1, S_1 / \nu_0) \\
z_0 \sim \Normal(0,  R / N)\\
S_1 \text{ is the Schur's complement of } R_{11} \text{ in } R
\end{cases}
\end{equation}
Note that this induces:
$$S_1 \sim \mathcal{IW}(\nu_0 \bar S_1, \nu_0 + J + 1).$$
Our goal is to show that this generative distribution of $(z_0, z_1)$ approximates:
\begin{equation}\label{eq:generative-z0-z1-by-omega}
\begin{cases}
z_1 \mid \omega \sim \Normal\left(\mu_1,  \dfrac{\bar S_1}{\omega \nu_0}\right) \\
z_0 \mid \omega \sim \Normal \left(0,  \dfrac{\bar R}{\omega N}\right)\\
\omega \sim \Gam((\nu_0 + 2)/ 2, \nu_0 / 2).
\end{cases}
\end{equation}

\subsection*{Approach 0}

We can easily show that the marginal distribution of $z_0$ is the same in both models. Can we show that $p(z_1 \mid z_0)$ from~\eqref{eq:generative-z0-z1-by-R} approximates $\tilde{p}(z_1 \mid z_0)$ from~\eqref{eq:generative-z0-z1-by-omega}?

Write
\[
z_0=
\begin{pmatrix}
z_{01}\\ z_{0,-1}
\end{pmatrix},
\qquad
\delta=z_{0,-1}-\mu_1z_{01},
\qquad
\kappa=\nu_0+Nz_{01}^2.
\]
Under the inverse-Wishart model,
\[
R\mid z_0
\sim
\mathcal{IW}\left(
\nu_0\Rbar+Nz_0z_0^\top,\nu_0+J+2
\right).
\]
Using the block inverse-Wishart identities,
\[
z_1\mid S_1,z_0
\sim
\Normal\left(
\mu_1+\frac{Nz_{01}}{\kappa}\delta,
\frac{S_1}{\kappa}
\right),
\]
and
\[
S_1\mid z_0
\sim
\mathcal{IW}\left(
\nu_0\bar S_1+
\frac{\nu_0 N}{\kappa}\delta\delta^\top,
\nu_0+J+2
\right).
\]
Thus the exact conditional distribution \(p(z_1\mid z_0)\) is centered at
\[
\mu_1+\frac{Nz_{01}}{\nu_0+Nz_{01}^2}
(z_{0,-1}-\mu_1z_{01}),
\]
not at \(\mu_1\). It also contains a rank-one update to the Schur complement
scale through
\[
\frac{\nu_0 N}{\kappa}\delta\delta^\top.
\]

In the shared-\(\omega\) approximation,
\[
z_0\mid\omega
\sim
\Normal\left(0,\frac{\bar R}{N\omega}\right),
\qquad
z_1\mid\omega
\sim
\Normal\left(\mu_1,\frac{\bar S_1}{\nu_0\omega}\right),
\]
and \(z_0\) and \(z_1\) are conditionally independent given \(\omega\).
Conditioning on \(z_0\) updates \(\omega\), but it does not shift the
conditional mean of \(z_1\). Therefore
\[
\E_{\mathrm{shared}\text{-}\omega}(z_1\mid z_0)=\mu_1.
\]

This means that the shared-\(\omega\) approximation does not generally match
the exact conditional distribution \(p(z_1\mid z_0)\). It ignores the
\(z_0\)-dependent mean shift
\[
\frac{Nz_{01}}{\nu_0+Nz_{01}^2}
(z_{0,-1}-\mu_1z_{01})
\]
and the rank-one covariance update
\[
\frac{\nu_0 N}{\nu_0+Nz_{01}^2}
(z_{0,-1}-\mu_1z_{01})(z_{0,-1}-\mu_1z_{01})^\top.
\]
The approximation is most plausible when this conditional residual effect is
small, for example when \(\nu_0\) is large, \(z_{01}\) is small, or
\(z_{0,-1}-\mu_1z_{01}\) is small. This is a useful baseline comparison to
return to later if we want to bound the approximation error more carefully.

\subsection{Strategy 1}

Strategy 1 starts from the normal--Gamma representation of the marginal prior
for the LD column:
\[
z_1\mid \omega
\sim
\Normal_{J-1}\left(
\mu_1,\frac{\bar S_1}{\nu_0\omega}
\right),
\qquad
\omega\sim
\Gam\left(\frac{\nu_0+3}{2},\frac{\nu_0}{2}\right).
\]
The question is whether, after introducing this \(\omega\), the noise vector
\(z_0\) can be approximated by
\[
z_0\mid \omega
\approx
\Normal_J\left(0,\frac{\Rbar}{N\omega}\right).
\]

For clarity, work with the \(d_1=1\) version. Conditional on the LD column
\(z_1\) and Schur complement \(S_1\), the covariance matrix is
\[
R=
\begin{pmatrix}
1 & z_1^\top\\
z_1 & S_1+z_1z_1^\top
\end{pmatrix}.
\]
Thus
\[
z_0\mid z_1,S_1
\sim
\Normal_J\left(
0,
\frac1N
\begin{pmatrix}
1 & z_1^\top\\
z_1 & S_1+z_1z_1^\top
\end{pmatrix}
\right).
\]
Equivalently,
\[
z_{01}\sim\Normal(0,1/N),
\qquad
z_{0,-1}\mid z_{01},z_1,S_1
\sim
\Normal_{J-1}\left(z_{01}z_1,\frac{S_1}{N}\right).
\]

The column mixture above can be viewed as replacing the random Schur complement
by the scalar approximation
\[
S_1\approx \frac{\bar S_1}{\omega}.
\]
Under this approximation,
\[
z_0\mid z_1,\omega
\approx
\Normal_J\left(
0,
\frac1N
\begin{pmatrix}
1 & z_1^\top\\
z_1 & z_1z_1^\top+\bar S_1/\omega
\end{pmatrix}
\right).
\]
Equivalently,
\[
z_{01}\sim\Normal(0,1/N),
\qquad
z_{0,-1}\mid z_{01},z_1,\omega
\sim
\Normal_{J-1}\left(
z_{01}z_1,\frac{\bar S_1}{N\omega}
\right).
\]

This is the direct consequence of Strategy 1. It is not the same as
\[
\Normal_J\left(0,\frac{\Rbar}{N\omega}\right),
\qquad
\Rbar=
\begin{pmatrix}
1 & \mu_1^\top\\
\mu_1 & \bar S_1+\mu_1\mu_1^\top
\end{pmatrix}.
\]
The direct Strategy 1 covariance is
\[
\frac1N
\begin{pmatrix}
1 & z_1^\top\\
z_1 & z_1z_1^\top+\bar S_1/\omega
\end{pmatrix},
\]
whereas the desired shared-\(\omega\) covariance is
\[
\frac1{N\omega}
\begin{pmatrix}
1 & \mu_1^\top\\
\mu_1 & \bar S_1+\mu_1\mu_1^\top
\end{pmatrix}.
\]
Therefore, starting from the \(z_1\) mixture does not by itself justify scaling
the entire noise covariance by \(1/\omega\). It justifies scaling the
conditional Schur-complement noise \(S_1/N\) by \(1/\omega\).

To reach the desired shared-\(\omega\) likelihood, we would need additional
approximations, for example:
\begin{enumerate}[leftmargin=*]
\item replace \(z_1\) by its center \(\mu_1\) in the noise covariance;
\item also scale the active-coordinate and rank-one parts of the covariance by
\(1/\omega\).
\end{enumerate}
The first approximation is natural when LD-column uncertainty is small. The
second is less directly implied by the column mixture. It may still be useful
computationally, but it should be treated as an extra approximation rather
than a consequence of the SER-IW block decomposition.

This suggests an alternative Strategy 1 approximation:
\[
z_0\mid z_1,\omega
\approx
\Normal_J\left(
0,
\frac1N
\begin{pmatrix}
1 & z_1^\top\\
z_1 & z_1z_1^\top+\bar S_1/\omega
\end{pmatrix}
\right),
\]
or a stabilized version replacing \(z_1\) by \(\mu_1\) only in the covariance.
This is more faithful to the column-first derivation, but it is less convenient
than the current shared-\(\omega\) model because the likelihood covariance
depends on the active column.


\subsection{Strategy 2}

Strategy 2 starts from the normal--Gamma representation of the marginal noise
vector:
\[
z_0\mid \omega
\sim
\Normal_J\left(0,\frac{\Rbar}{N\omega}\right),
\qquad
\omega\sim
\Gam\left(\frac{\nu_0+2}{2},\frac{\nu_0}{2}\right).
\]
The question is whether, after introducing this \(\omega\), the LD column can
be approximated by
\[
z_1\mid \omega
\approx
\Normal_{J-1}\left(
\mu_1,\frac{\bar S_1}{\nu_0\omega}
\right).
\]

Using the block form
\[
\Rbar=
\begin{pmatrix}
1 & \mu_1^\top\\
\mu_1 & \bar S_1+\mu_1\mu_1^\top
\end{pmatrix},
\]
the noise mixture implies
\[
z_{01}\mid \omega
\sim
\Normal\left(0,\frac{1}{N\omega}\right),
\]
and
\[
z_{0,-1}\mid z_{01},\omega
\sim
\Normal_{J-1}\left(
\mu_1z_{01},
\frac{\bar S_1}{N\omega}
\right).
\]
Equivalently, the conditional noise residual
\[
\delta_0=z_{0,-1}-\mu_1z_{01}
\]
satisfies
\[
\delta_0\mid \omega
\sim
\Normal_{J-1}\left(
0,\frac{\bar S_1}{N\omega}
\right),
\qquad
\delta_0\perp z_{01}\mid\omega.
\]

Now use the observation equation
\[
x=
\beta
\begin{pmatrix}
1\\ z_1
\end{pmatrix}
+z_0.
\]
For \(\gamma=1\),
\[
x_1=\beta+z_{01},
\qquad
x_{-1}=\beta z_1+z_{0,-1}.
\]
Subtracting \(\mu_1x_1\) gives
\[
x_{-1}-\mu_1x_1
=
\beta(z_1-\mu_1)+\delta_0.
\]
Thus, conditional on \(\omega\),
\[
x_{-1}-\mu_1x_1
\mid \beta,z_1,\omega
\sim
\Normal_{J-1}\left(
\beta(z_1-\mu_1),
\frac{\bar S_1}{N\omega}
\right).
\]

If we now place the working prior
\[
z_1\mid\omega
\sim
\Normal_{J-1}\left(
\mu_1,\frac{\bar S_1}{\nu_0\omega}
\right),
\]
then the conditional residual has the marginal distribution
\[
x_{-1}-\mu_1x_1
\mid \beta,\omega
\sim
\Normal_{J-1}\left(
0,
\frac{1}{\omega}
\left(\frac1N+\frac{\beta^2}{\nu_0}\right)\bar S_1
\right).
\]
This is close in form to the exact model-B conditional distribution, except
that \(x_1^2\) is replaced by \(\beta^2\):
\[
x_{-1}-\mu_1x_1
\mid x_1
\quad\text{exactly has scale roughly}\quad
\left(\frac1N+\frac{x_1^2}{\nu_0}\right)\bar S_1.
\]
Therefore Strategy 2 gives a coherent full-vector Gaussian model, but it
changes the conditioning logic: the nonactive residual is explained through
\(\beta(z_1-\mu_1)\), rather than through \(x_1(z_1-\mu_1)\).

Compared with Strategy 1, Strategy 2 more naturally justifies scaling the
entire observation-noise covariance by \(1/\omega\), because this is how the
strategy starts. The additional approximation is instead in the LD-column
prior: we attach the same noise scale \(\omega\) to \(z_1-\mu_1\), even though
the exact inverse-Wishart dependence between \(z_0\) and \(z_1\) is more
complicated.

This strategy leads directly to the shared-\(\omega\) model:
\[
z_0\mid\omega
\sim
\Normal_J\left(0,\frac{\Rbar}{N\omega}\right),
\qquad
z_1\mid\omega
\sim
\Normal_{J-1}\left(
\mu_1,\frac{\bar S_1}{\nu_0\omega}
\right),
\]
\[
x=\beta
\begin{pmatrix}
1\\z_1
\end{pmatrix}
+z_0.
\]
In this sense, Strategy 2 is a more direct justification for the current
shared-\(\omega\) likelihood, while Strategy 1 is a more direct justification
for scaling the Schur-complement part of the LD uncertainty.


\section{SER Variational Inference for Shared-\texorpdfstring{\(\omega\)}{omega}}

Use the mean-field approximation
\[
q(\gamma=j,\beta,z_j,\omega_j)
=
\pi_jq_j(\beta)q_j(z_j)q_j(\omega_j),
\]
with
\[
q_j(\beta)=\Normal(m_{\beta j},v_{\beta j}),
\qquad
q_j(z_j)=\Normal(m_{zj},V_{zj}),
\qquad
q_j(\omega_j)=\Gam(\alpha_{\omega j},\beta_{\omega j}).
\]
Let
\[
\bar\omega_j=\E_j(\omega_j),
\qquad
\bar\beta_j=m_{\beta j},
\qquad
\overline{\beta^2}_j=v_{\beta j}+m_{\beta j}^2.
\]
Define
\[
\Om=\Rbar^{-1},
\qquad
C_j(x)=x^\top\Om x-x_j^2.
\]
For the current \(q_j(z_j)\), define
\begin{align*}
Q_j
&=
\tr(\bar S_j^{-1}V_{zj})
+(m_{zj}-\mu_j)^\top\bar S_j^{-1}(m_{zj}-\mu_j),\\
A_j&=1+Q_j,\\
B_j
&=
x_j+(m_{zj}-\mu_j)^\top\bar S_j^{-1}
(x_{-j}-\mu_jx_j).
\end{align*}
The beta update is
\[
v_{\beta j}
=
\left(\sigma_0^{-2}+N\bar\omega_jA_j\right)^{-1},
\qquad
m_{\beta j}
=
v_{\beta j}N\bar\omega_jB_j.
\]
The relaxed-column update is
\[
V_{zj}
=
\frac{\bar S_j}
{\bar\omega_j(N\overline{\beta^2}_j+\nu_0)},
\]
\[
m_{zj}
=
\mu_j+
\frac{N\bar\beta_j}
{N\overline{\beta^2}_j+\nu_0}
(x_{-j}-\mu_jx_j).
\]
The Gamma update is
\[
\alpha_{\omega j}
=
a_\omega+\frac{J}{2}+\frac{J-1}{2},
\]
\[
\beta_{\omega j}
=
b_\omega
+\frac12
\left[
N\left\{
x^\top\Om x-2\bar\beta_jB_j+\overline{\beta^2}_jA_j
\right\}
+\nu_0Q_j
\right].
\]
The coordinate weights are
\[
\pi_j
=
\frac{\exp(\mathcal L_j)}
{\sum_{k=1}^J\exp(\mathcal L_k)},
\]
where \(\mathcal L_j\) is the local variational lower bound including the
likelihood, the priors for \(\beta,z_j,\omega_j\), and the entropies of
\(q_j(\beta),q_j(z_j),q_j(\omega_j)\).

\section{PIP Distortion View}

The ideal SER weight has the form
\[
\ell_j^{\mathrm{ideal}}
=
\mathrm{const}
+
\frac{N}{2}
\frac{N\sigma_0^2}{1+N\sigma_0^2}
x_j^2.
\]
The shared-\(\omega\) approximation produces some local weight
\[
\ell_j^{\mathrm{sh}\text{-}\omega}
=
\ell_j^{\mathrm{ideal}}+\delta_j.
\]
The goal is not necessarily to force \(\delta_j=0\). Instead, the useful target
is to understand and bound the spread
\[
\Delta
=
\max_j\delta_j-\min_j\delta_j.
\]
If \(\Delta\) is small, then the approximate PIPs are uniformly close to the
ideal PIPs in the sense that their odds are distorted by at most a factor
\(\exp(\Delta)\). This gives a concrete way to judge whether the
shared-\(\omega\) approximation is acceptable.

\end{document}
